\chapter{Transport System Implementation}  \label{transportSystemImpl}

\section{Candidate Gathering}
CTaps follows the recommended procedure for Candidate Gathering from RFC 9623
very closely. CTaps represents the candidates in a tree, using a GNode struct from
GLib to represent a N-ary tree. CTaps first builds the entire tree, resolving
all possible local interfaces, all protocols and all possible resolutions of
remote hostnames respectively. If the client has specified multiple remote endpoints,
then they are treated the same way as each invidvidual resolution from a hostname
endpoint. After the entire tree has been built, CTaps prunes the branches which
are not compatible with the required or prohibited selection properties. After pruning,
the tree is sorted according to the stated selection properties, as described in RFC 9623 
\cite{brunstromImplementingInterfacesTransport2025}.

\section{Candidate Racing}

\section{Protocol Implementations}
In order to make the CTaps design as modular as possible and ease the
process of adding more protocols, CTaps uses an internal interface for a protocol,
which abstracts away the implementation details. This interface consists of
a struct containing the name of the protocols, and function pointers for
the necessary functions for the wider CTaps library to interact with the protocol.
Each protocol implementation is this way kept as small as possible, and is
isolated from the rest of the library.
 
\begin{figure}[H]
\begin{minted}{C}
typedef struct ProtocolImplementation {
  const char* name;
  SelectionProperties selection_properties;
  int (*init)(struct Connection* connection, InitDoneCb init_done_cb);
  int (*send)(struct Connection*, Message*);
  int (*receive)(struct Connection*, ReceiveMessageRequest receive_cb);
  int (*listen)(struct SocketManager* socket_manager);
  int (*stop_listen)(struct SocketManager*);
  int (*close)(const struct Connection*);
} ProtocolImplementation;
\end{minted}
\caption{Connection typedef struct}
\end{figure} \todo{Update this when having finalized the code}

When implementing a new protocol, all that is needed is an
implementation of the struct, filling out the function pointers
with custom behaviour for the given protocol:

\begin{figure}[H]
\begin{minted}{C}
int udp_init(Connection* connection, InitDoneCb init_done_cb);
int udp_close(const Connection* connection);
int udp_send(Connection* connection, Message* message);
int udp_receive(Connection* connection, ReceiveMessageRequest receive_message_cb);
int udp_listen(struct SocketManager* socket_manager);
int udp_stop_listen(struct SocketManager* listener);

static ProtocolImplementation udp_protocol_interface = {
    .name = "UDP",
    .selection_properties = {
      .selection_property = {
        get_selection_property_list(create_property_initializer)
        [RELIABILITY] = {.value = {.simple_preference = PROHIBIT}},
      }
    },
    .send = udp_send,
    .init = udp_init,
    .receive = udp_receive,
    .close = udp_close,
    .listen = udp_listen,
    .stop_listen = udp_stop_listen
};
\end{minted}
\caption{UDP protocol implementation struct}
\end{figure} \todo{Update this when having finalized the code}
 
These implementations are kept track of in a central array,
and new protocols can be dynamically registered by consumers
of the library. Implementations of the interface also register
a list of the applicable selection properties, allowing the library
to later prune and sort the Candidate Tree protocols, when selecting
candidates.
